\documentclass{article}
\usepackage{amsmath,amssymb}
\usepackage[vocab={realfunctions,optimization,calculus}]{leanify}
\newtheorem{lemma}{Lemma}
\begin{document}

\section{Notation}

We write $D^n f$ for the $n$th derivative of $f$.
We write $S(f)$ for the second derivative of $f$.
We write $T_n(f,a)$ for the Taylor polynomial of degree $n$ of $f$ at $a$.
We write $\ell(f,a,b)$ for the arc length of the graph of $f$ on $[a,b]$.
We write $\mu(f,a,b)$ for the average value of $f$ on $[a,b]$.
We write $A(f,g,a,b)$ for the area between the graphs of $f$ and $g$ on $[a,b]$.
We write $V(f,a,b)$ for the volume of the solid of revolution of $f$ on $[a,b]$.
For every real number $x$, let $p(x)\coloneqq x^3$.

\section{Derivatives}

\begin{lemma}\label{second}
Let $f$ be a twice differentiable function. Then $S(f)(x)=D^2f(x)$ for every real number $x$.
\end{lemma}

\begin{lemma}\label{times}
Every twice continuously differentiable function is twice differentiable.
\end{lemma}

\begin{lemma}\label{taylor}
Let $f$ be a $3$ times differentiable function and let $a$ be a real number. Then $T_0(f,a)(x)=f(a)$ for every real number $x$.
\end{lemma}

\begin{lemma}\label{critical}
Let $f$ be a differentiable function and let $x$ be a real number. If $x$ is a local minimizer of $f$, then $x$ is a critical point of $f$.
\end{lemma}

\begin{lemma}\label{inflection}
$0$ is an inflection point of $p$.
\end{lemma}

\begin{lemma}\label{antiderivative}
Let $f$ and $F$ be functions. If $F$ is an antiderivative of $f$, then $F$ is differentiable.
\end{lemma}

\section{Integrals}

\begin{lemma}\label{length}
Let $f$ be a continuously differentiable function and let $a$ and $b$ be real numbers with $a\le b$. Then $\ell(f,a,b)\ge b-a$.
\end{lemma}

\begin{lemma}\label{average}
Let $f$ be a continuous function and let $a$ and $b$ be real numbers with $a<b$. Then there is $c\in[a,b]$ such that $f(c)=\mu(f,a,b)$.
\end{lemma}

\begin{lemma}\label{area}
Let $f$ and $g$ be continuous functions and let $a$ and $b$ be real numbers with $a\le b$. Then $A(f,g,a,b)=A(g,f,a,b)$ and $V(f,a,b)\ge0$.
\end{lemma}

\begin{lemma}\label{words-without-noun}
Let $f$ be twice differentiable with $f^{\prime\prime}(x)\ge0$ for all $x$. Then $f$ is convex.
\end{lemma}

\end{document}
